\section{Method}
\label{sec:method}

Motivated by the findings in
Section~\ref{sec:empirical-self-speculate:findings}, we train a single
tool-calling LLM agent $\pi_\theta$ to act both as the task-solving agent and as
its own next-call speculator. We first define the
agent task reward and the speculation reward used to score predicted tool calls.
Second, we describe the joint agent--speculator RL method, which turns
current agent trajectories into on-policy speculation targets. Finally, we describe 
the training recipe used to stabilize the shared policy for both modes.

\subsection{Task and Reward Formulation}
\label{sec:method:task-formulation}

Given a user query $x$, the agent interacts with an environment through a
sequence of reasoning steps, tool calls, and tool observations. We denote an
agent trajectory by
$\tau=(x,r_1,a_1,o_1,\ldots,r_T,a_T,o_T,y)$, where $r_t$ is the reasoning
before the $t$-th tool call, $a_t$ is the tool call, $o_t$ is the returned
observation, and $y$ is the final answer. Each tool call is a structured action
$a_t=(n_t,\sigma_t)$, with tool name $n_t$ and argument dictionary $\sigma_t$.
The agent reward $R_{\mathrm{ag}}(\tau)$ is evaluated on the final answer
$y$ against the ground-truth answer $y^*$ defined in the environment.

The speculation task is derived from the same trajectory and is implemented by
the same policy $\pi_\theta$. Agent and speculator modes therefore share model
parameters, with the speculator mode triggered simply by appending a fixed speculation suffix 
to the input context. For each tool-call turn $t$, we take the prefix immediately before
the call,
$h_t=(x,r_1,a_1,o_1,\ldots,a_{t-1})$, append a fixed speculation suffix
$s_{\mathrm{sp}}$, and ask the policy to emit the next structured tool
call, $\hat a_t\sim\pi_\theta(\cdot\mid h_t\oplus s_{\mathrm{sp}})$. The target
is the agent's own next call $a_t$ from the sampled trajectory. Since $a_t$
comes from the current policy, the targets change as the agent changes, keeping
speculation training on-policy with the deployed agent.

We score a predicted call $\hat a_t=(\hat n_t,\hat\sigma_t)$ against the target
call $a_t=(n_t,\sigma_t)$ using a name-conditioned argument reward. The tool-name
score is $R_{\mathrm{name}}=\mathbf{1}[\hat n_t=n_t]$. If the tool name matches,
we compare argument dictionaries with macro token-F1 over target argument keys:
\[
R_{\mathrm{args}}
= \frac{1}{|\mathrm{keys}(\sigma_t)|}
\sum_{k\in\mathrm{keys}(\sigma_t)}
\mathrm{F1}\!\left(\hat\sigma_t[k],\sigma_t[k]\right).
\]
Missing predicted arguments receive score zero, and extra predicted arguments
do not receive credit. Argument values are normalized to strings before
token-F1. For tools with no arguments, we set $R_{\mathrm{args}}=1$ after a
correct tool-name match. The final speculation reward is
\begin{equation}
R_{\mathrm{sp}} = R_{\mathrm{name}}\cdot R_{\mathrm{args}}.
\label{eq:reward}
\end{equation}
This construction matches the exact-match reuse requirement for speculation: a wrong tool name cannot
reuse the pre-executed result, while a correct tool name can still receive
partial credit for close argument values. Macro averaging gives each target
argument key equal weight, avoiding pooled token-F1 behavior where long argument
values dominate shorter but equally important keys.

\subsection{Joint Agent--Speculator RL}
\label{sec:method:dual-mode-rollout}

Agent and speculator samples differ substantially. The agent mode generates long,
reasoning-heavy trajectories, whereas the speculator mode generates a single structured tool call 
from a partial trajectory. Their different output distributions and rewards make direct joint training difficult.
We therefore keep the data collection coupled but the policy updates separated. 
Each iteration constructs fresh on-policy data for both modes, while each optimizer step uses only the objective selected for
the current training iteration.

At RL iteration $k$, we sample a query batch $Q\subset\mathcal{D}$ and generate
$G_{\mathrm{ag}}$ agent trajectories for each $x\in Q$. These trajectories serve as standard agent rollouts: their final rewards
$R_{\mathrm{ag}}$ are normalized within each query group to produce
DAPO-style relative advantages for the agent objective~\citep{yu2025dapo}. At
the same time, they provide the speculation supervision induced by the current
policy. For each realized tool-call turn, we take the prefix $h_t$ and the
actual next call $a_t$ from the sampled trajectory, form the speculator input
$h_t\oplus s_{\mathrm{sp}}$, and treat $a_t$ as the target call.

A fixed schedule $m(k)\in\{\mathrm{ag},\mathrm{sp}\}$ selects which objective is
optimized at iteration $k$. If $m(k)=\mathrm{ag}$, we update
$\pi_\theta$ with the agent rollouts and advantages from $R_{\mathrm{ag}}$.
If $m(k)=\mathrm{sp}$, we update the same $\pi_\theta$ with the speculator
rollouts and advantages from $R_{\mathrm{sp}}$. Both modes use the same
DAPO policy optimization~\citep{yu2025dapo}; they differ in the input context, rollout structure,
and reward function. In our final runs, $m(k)$ follows a repeated 4:8
pattern: four consecutive agent updates followed by eight consecutive
speculator updates. Algorithm~\ref{alg:training-recipe} summarizes the full procedure.

\subsection{Stabilizing Dual-Mode RL}
\label{sec:method:training-recipe}

Training a shared policy for both task solving and self-speculation is
non-trivial. The two modes produce different output distributions and optimize
different rewards, so direct joint training can suffer from mode interference or
collapse into unproductive tool-calling behavior. We therefore use a training
recipe with three components. \ding{182} Before RL, we run a short SFT warmup
on a mixture of successful agent trajectories and their corresponding
intermediate-trajectory speculation examples. The agent examples supervise the
full successful trajectory, while each speculation example takes a partial
trajectory, appends $s_{\mathrm{sp}}$, and supervises the model to produce the
ground-truth next tool call. This gives the model well-formed behavior in both
modes before RL begins. \ding{183} We reset optimizer state whenever the
schedule switches between agent and speculator updates. Since the two modes optimize different rewards, 
momentum and adaptive statistics differ between the two modes, which can strongly affect subsequent optimization
steps~\citep{sutskever2013momentum,kingma2015adam}; the reset reduces carryover
from one mode's objective to the next. \ding{184} The agent reward includes a
penalty for repeated meaningless tool calls, discouraging a common failure mode
in which the agent continues calling tools without making task progress.

\begin{algorithm}[t]
\small
\begin{algorithmic}[1]
\Require Policy $\pi_\theta$, query set $\mathcal{D}$, speculation suffix $s_{\mathrm{sp}}$, alternation schedule $m(k)$, group sizes $G_{\mathrm{ag}},G_{\mathrm{sp}}$
\State $m_{\mathrm{prev}}\gets \varnothing$
\For{RL iteration $k = 1,2,\ldots$}
    \State Sample query batch $Q \sim \mathcal{D}$
    \State Generate agent groups $\mathcal{T}=\{\tau_i(x)\}_{x\in Q,\,i=1}^{G_{\mathrm{ag}}}$
    \State Build speculation pairs $\mathcal{S}=\{(h_t,a_t)\!: t\in\tau,\tau\in\mathcal{T}\}$
    \State Generate speculator groups $\widehat{\mathcal{A}}=\{\hat a_{t,j}\sim\pi_\theta(\cdot\mid h_t\oplus s_{\mathrm{sp}})\}_{(h_t,a_t)\in\mathcal{S},\,j=1}^{G_{\mathrm{sp}}}$
    \If{$m(k)\ne m_{\mathrm{prev}}$}
        \State Reset optimizer state
    \EndIf
    \If{$m(k)=\mathrm{ag}$}
        \State DAPO update $\pi_\theta$ on grouped agent rewards $R_{\mathrm{ag}}(\mathcal{T})$
    \Else
        \State DAPO update $\pi_\theta$ on grouped speculation rewards $\{R_{\mathrm{sp}}(\hat a_{t,j},a_t)\}_{(h_t,a_t)\in\mathcal{S},\,j=1}^{G_{\mathrm{sp}}}$
    \EndIf
    \State $m_{\mathrm{prev}}\gets m(k)$
\EndFor
\end{algorithmic}
\caption{Joint Agent-Speculator RL}
\label{alg:training-recipe}
\end{algorithm}

